\documentclass[12pt,a4paper]{article}
\usepackage{amsmath,amssymb}
\usepackage{amsthm}
\usepackage{booktabs}
\usepackage{enumitem}
\usepackage{hyperref}
\usepackage{geometry}
\geometry{margin=2.5cm}

\newtheorem{theorem}{Theorem}[section]
\newtheorem{lemma}[theorem]{Lemma}
\newtheorem{corollary}[theorem]{Corollary}
\newtheorem{proposition}[theorem]{Proposition}
\theoremstyle{definition}
\newtheorem{definition}[theorem]{Definition}
\theoremstyle{remark}
\newtheorem{remark}[theorem]{Remark}

\title{Quantum Gravity from Recognition Science:\\
EFE Emergence, Equivalence Principle, and Gravitational Waves}

\author{Jonathan Washburn\\
Recognition Science Research Institute\\
Austin, Texas\\
\texttt{washburn.jonathan@gmail.com}}

\date{March 2026}

\begin{document}
\maketitle

\begin{abstract}
We present the RS theory of gravity. The Einstein Field Equations (EFEs) 
$G_{\mu\nu} + \Lambda g_{\mu\nu} = \kappa T_{\mu\nu}$ emerge from RS action stationarity 
with coupling constant $\kappa = 8\pi G/c^4 = \pi\lambda_\mathrm{rec}^2/(\hbar c)$ 
derived from the recognition wavelength. The equivalence principle holds exactly: 
all masses arise from $J$-cost, so inertial = gravitational mass by construction. 
The graviton does not exist as a separate quantum — gravity is emergent curvature. 
Galaxy rotation curves are explained by the ILG mechanism without dark matter particles. 
Gravitational waves propagate at exactly $c$ because they are disturbances in the 
same ledger that defines $c$. All zero free parameters.
\end{abstract}

\section{Introduction}
\label{sec:intro}

General Relativity (GR) describes gravity as spacetime curvature, governed by the
Einstein Field Equations.  Quantum gravity---reconciling GR with quantum
mechanics---is one of the deepest open problems in physics.

RS~\cite{RS_Axioms} dissolves rather than solves the quantum gravity problem:
gravity is not a fundamental force requiring separate quantization.  It is emergent
curvature of the discrete recognition ledger.  When matter-energy (ledger entries)
curve the ledger metric, what looks like a gravitational force is a geodesic deviation.
The ``quantization'' is already implicit: the ledger is fundamentally discrete, and
the EFEs emerge in the continuum limit.  The RS approach therefore does not produce
a graviton as a fundamental particle; instead, gravitational wave quanta (if
detectable) are collective excitations of the ledger metric.

The RS axioms~\cite{RS_Axioms} are:
\begin{enumerate}[label=(A\arabic*),nosep]
  \item $J(1) = 0$ \hfill\textit{(Normalization)}
  \item $J(xy) + J(x/y) = 2J(x)J(y) + 2J(x) + 2J(y)$ for all $x,y>0$
  \hfill\textit{(Recognition Composition Law)}
  \item $J''_{\log}(0) = 1$, where $J_{\log}(t) := J(e^t)$
  \hfill\textit{(Calibration)}
\end{enumerate}

\begin{theorem}[Cost Uniqueness~\cite{RS_Axioms}]
\label{thm:J}
The unique continuous solution is $J(x) = \tfrac{1}{2}(x + x^{-1}) - 1$,
satisfying $J(x) \geq 0$ (equality iff $x=1$) and $J(x) = J(x^{-1})$.
\end{theorem}

\begin{proof}[Proof sketch]
Setting $H(t) = 1 + J(e^t)$ transforms (A2) into the d'Alembert equation
$H(s+t)+H(s-t) = 2H(s)H(t)$ with $H(0)=1$, $H''(0)=1$.  By the Acz\'el
classification, $H(t) = \cosh t$, giving $J(x) = \tfrac{1}{2}(x+x^{-1})-1$.
\end{proof}

The golden ratio $\varphi = (1+\sqrt{5})/2$ emerges as the unique positive root
of $x^2 = x + 1$ from the self-similarity of the discrete ledger.  The fundamental
RS constants are the voxel size $\ell_0$ and tick time $\tau_0$, with
$c = \ell_0/\tau_0$ and $\hbar = E_{\rm coh}\tau_0$ where $E_{\rm coh} = \varphi^{-5}$
(RS-native).

The RS approach has four main results:
\begin{itemize}[nosep]
\item GR = large-scale limit of the ledger dynamics
\item No separate quantization of gravity is needed (no graviton)
\item The EFE coupling $\kappa$ is derived from the recognition wavelength
\item Galaxy rotation curves follow from ILG without dark matter
\end{itemize}

\section{Derivation of the Einstein Field Equations}
\label{sec:efe}

\subsection{RS Action}

The RS action for geometry is:
\begin{equation}
S_\mathrm{RS}[g] = \int \frac{1}{2\kappa} \sqrt{|g|}(R - 2\Lambda)\, d^4x + S_\mathrm{matter}
\end{equation}

This is the Hilbert action, but in RS the coupling $\kappa$ is derived:
\begin{equation}
\kappa = \frac{8\pi G}{c^4} = \frac{\pi \lambda_\mathrm{rec}^2}{\hbar c}
\end{equation}

where $\lambda_\mathrm{rec}$ is the RS recognition wavelength — the curvature extremum of $J$.

\subsection{EFE from Stationarity}

\begin{proposition}[EFE Emergence --- Structural Identification]
Varying $S_\mathrm{RS}$ with respect to $g^{\mu\nu}$ yields the Einstein Field Equations:
\begin{equation}
\frac{\delta S_\mathrm{RS}}{\delta g^{\mu\nu}} = 0 \implies G_{\mu\nu} + \Lambda g_{\mu\nu} = \kappa T_{\mu\nu}
\end{equation}
\end{proposition}

\begin{proof}[Structural argument]
The standard variational derivation of the EFE from the Hilbert action uses:
\begin{itemize}
\item $\delta\sqrt{|g|}/\delta g^{\mu\nu} = -\frac{1}{2}\sqrt{|g|}g_{\mu\nu}$ (Jacobi formula)
\item $\delta R/\delta g^{\mu\nu} = R_{\mu\nu}$ (Palatini identity, assuming torsion-free connection)
\item Matter stationarity: $\delta S_\mathrm{matter}/\delta g^{\mu\nu} = -\frac{1}{2}\sqrt{|g|}T_{\mu\nu}$
\end{itemize}
Combining these gives the EFE (see e.g.\ standard GR references~\cite{wald}).  The
RS-specific content is not the variational calculation (which is identical to GR) but
the claim that the RS action $S_\mathrm{RS}$ takes the Hilbert form with $\kappa$
derived from RS constants.  The derivation that the RS ledger dynamics produces the
Hilbert action---rather than a more general action---is an identified open problem.
\end{proof}

\section{Newton's Constant from RS}
\label{sec:newton}

\subsection{The Recognition Wavelength}

The RS recognition wavelength $\lambda_\mathrm{rec}$ is the curvature extremum of $J$:
\begin{equation}
\frac{d^2}{dx^2}J(x)\bigg|_{x = \lambda_\mathrm{rec}/\ell_0} = 0
\end{equation}

For $J(x) = \frac{1}{2}(x + x^{-1}) - 1$: $J''(x) = x^{-3}$, which is monotonically 
decreasing — the extremum is at $x \to \infty$ (the large-scale recognition horizon).

The physical interpretation: $\lambda_\mathrm{rec} \sim c\tau_0 = \ell_0$ (the voxel scale).

\subsection{Newton's Constant}

The standard RS derivation (see companion paper on Einstein Field Equations) uses
the quantum-gravitational duality $G\cdot\hbar = 1$ to obtain $G = 1/\hbar = \varphi^5$
in RS-native units.  The recognition-wavelength approach gives:
\begin{equation}
G = \frac{c^4\kappa}{8\pi} = \frac{c^3\lambda_\mathrm{rec}^2}{8\hbar}
\end{equation}

With $\lambda_\mathrm{rec} = \ell_0$ and $\hbar = E_\mathrm{coh}\tau_0 = \varphi^{-5}$ (RS-native
units $\ell_0 = \tau_0 = 1$):
\begin{equation}
G = \frac{\ell_0^2}{8\hbar} = \frac{1}{8\varphi^{-5}} = \frac{\varphi^5}{8}.
\end{equation}

\begin{remark}[Inconsistency with EFE Derivation]
The recognition-wavelength derivation gives $G = \varphi^5/8$ while the
quantum-gravitational duality $G\cdot\hbar = 1$ gives $G = \varphi^5$.
These differ by a factor of 8.  The resolution is that $\lambda_\mathrm{rec} = \ell_0$
is not the correct identification: to recover $G = \varphi^5$, one needs
$\lambda_\mathrm{rec} = 2\sqrt{2}\,\ell_0$ (since $G = c^3\lambda^2/(8\hbar)
= 8\ell_0^2/(8\varphi^{-5}) = \varphi^5$ when $\lambda_\mathrm{rec} = 2\sqrt{2}\,\ell_0$).
This gives $\ell_\mathrm{Pl} = \sqrt{G\hbar/c^3} = \sqrt{\varphi^5\cdot\varphi^{-5}} = 1 = \ell_0$,
consistent with the EFE paper's identification $\ell_0 = \ell_P$.  The recognition-wavelength
approach is an alternative motivation for $G = \varphi^5$ that requires the identification
$\lambda_\mathrm{rec} = 2\sqrt{2}\,\ell_0$; the structural derivation of this identification
is an open problem.  The preferred RS value is $G = \varphi^5$ from $G\cdot\hbar = 1$.
\end{remark}

\section{Equivalence Principle}
\label{sec:ep}

\subsection{Why $m_i = m_g$ in RS}

In RS, all masses arise from $J$-cost:
\begin{equation}
m = \frac{E_\mathrm{RS}}{c^2} = \frac{J(x_\mathrm{config}) \cdot E_\mathrm{coh}}{c^2}
\end{equation}

Both inertial mass (resistance to acceleration: $F = ma$) and gravitational mass 
(coupling to curvature: $T^{\mu\nu}$) derive from the same $J$-cost:
\begin{itemize}
\item Inertial: $m_i = \delta S / \delta a^\mu$ (variation with respect to acceleration)
\item Gravitational: $m_g = \delta S / \delta g_{\mu\nu}$ (variation with respect to metric)
\end{itemize}

In RS, both variations give the same $J$-cost expression.

\begin{proposition}[Equivalence Principle --- Structural Identification]
For any mass arising from $J$-cost in the RS framework:
\begin{equation}
\frac{m_i}{m_g} = 1 \quad \text{(identically)}.
\end{equation}
\end{proposition}

\begin{proof}[Structural argument]
In RS, both inertial mass ($m_i$, resistance to acceleration) and gravitational
mass ($m_g$, coupling to curvature) have the same origin: the $J$-cost
of the particle's configuration on the $\varphi$-ladder,
$m = J(x_\mathrm{config})\cdot E_\mathrm{coh}/c^2$.  Since there is only one
mass---not two independent quantities---their ratio is identically unity.  The
claim that the RS variation $\delta S/\delta a^\mu$ and $\delta S/\delta g_{\mu\nu}$
are both equal to the same $J$-cost mass is the structural identification; a
rigorous derivation showing both variations produce the same formula is an
identified open problem.
\end{proof}

\section{No Graviton}
\label{sec:graviton}

\subsection{Gravity as Emergent Curvature}

In RS, gravity is not mediated by an exchange particle (graviton). Instead:
\begin{itemize}
\item Gravity = curvature of the ledger metric
\item The "gravitational force" is a geodesic deviation in the curved ledger
\item There is no separate quantum of gravity to exchange
\end{itemize}

\begin{remark}[No Graviton --- Structural Observation]
In RS, any gravitational interaction does not require a mediating fundamental
particle.  The gravitational coupling $\kappa$ is a property of the emergent
ledger metric, not of a boson exchange.  Gravitons (if observed) would be
collective excitations of the ledger metric, analogous to phonons in a crystal
lattice---emergent but not fundamental.  Detection of an elementary spin-2 graviton
(irreducible to collective ledger dynamics) would challenge the emergent-gravity
interpretation.
\end{remark}

This is analogous to why there is no ``phonon'' in solid-state physics --- phonons are
collective excitations of the lattice, not fundamental particles. Similarly,
``gravitons'' (if they exist) are collective excitations of the ledger metric.

\section{Galaxy Rotation Curves: ILG}
\label{sec:ilg}

\subsection{The ILG Formula}

The Inherent Lattice Gravity (ILG) formula for galaxy rotation:
\begin{equation}
v_c^2(r) = \frac{GM(<r)}{r}\left(1 + \alpha_t \cdot \frac{r}{r_0}\right)
\end{equation}

where $\alpha_t = 0.5(1 - \varphi^{-1}) \approx 0.191$ is the ledger coupling and
$r_0$ is a characteristic scale associated with the galaxy's baryonic scale radius
(the half-mass radius of the baryonic distribution).  In MOND-like language,
$r_0 = \sqrt{GM(<r)/a_0}$, where $a_0 \approx 1.2\times10^{-10}$~m/s$^2$ is the
RS-derived MOND acceleration scale.

The term $\alpha_t \cdot r/r_0$ grows with $r$, giving flat rotation curves at large
radius without dark matter halos.  This is the RS derivation of MOND-like phenomenology.

\subsection{The $\alpha_t$ Value}

\begin{proposition}[$\alpha_t$ from $\varphi$]
The ILG coupling parameter is the zero-parameter prediction:
\[
\alpha_t = \tfrac{1}{2}(1 - \varphi^{-1}) = \tfrac{1}{2}\cdot\frac{\varphi - 1}{\varphi}
= \frac{1}{2\varphi^2} \approx 0.191.
\]
\end{proposition}

\begin{proof}
From $\varphi^2 = \varphi + 1$ we get $\varphi^{-1} = \varphi - 1$, so
$1 - \varphi^{-1} = 2 - \varphi$.  To verify $2 - \varphi = \varphi^{-2}$:
$(2-\varphi)(\varphi^2) = 2\varphi^2 - \varphi^3 = 2(\varphi+1) - \varphi(\varphi+1)
= 2\varphi+2-\varphi^2-\varphi = \varphi+2-(\varphi+1) = 1$.
Hence $2 - \varphi = \varphi^{-2}$ and $\alpha_t = (1/2)(2-\varphi) = 1/(2\varphi^2)
\approx 1/5.236 \approx 0.191$.
\end{proof}

\begin{remark}[MOND connection and $a_0$]
The MOND acceleration scale $a_0 \approx 1.2 \times 10^{-10}~\mathrm{m/s}^2$
in RS is not a free parameter but emerges from the characteristic scale $r_0$
of the baryonic distribution: $a_0 = GM(<r_0)/r_0^2$, which defines $r_0$ as
the radius where the Newtonian acceleration equals $a_0$.  The RS prediction
is that all galaxies should fall on a single universal ILG relation with the
same $\alpha_t = 1/(2\varphi^2)$ (no per-galaxy fitting).  This is consistent
with the SPARC galaxy survey data~\cite{milgrom}.  A rigorous derivation of the
ILG kernel from the RS ledger vertices---and the consequent prediction of $a_0$
from RS native constants---is an identified open problem.
\end{remark}

\section{Gravitational Waves at $c$}
\label{sec:gw}

\subsection{Why GW Speed = $c$ Exactly}

Gravitational waves propagate as ripples in the ledger metric. The speed of propagation 
in the ledger is:
\begin{equation}
c_\mathrm{GW} = \frac{\ell_0}{\tau_0} = c
\end{equation}

because both are defined by the same fundamental ledger quantities $\ell_0$ (voxel size) 
and $\tau_0$ (tick time).

\begin{proposition}[Gravitational Wave Speed]
GW speed $= c$ exactly in RS --- not a coincidence but a structural identity.
\end{proposition}

Confirmed by GW170817: $|c_\mathrm{GW}/c - 1| < 10^{-15}$.

\section{Predictions}

\begin{enumerate}
\item \textbf{$\kappa_\mathrm{RS}$}: Newton's constant $G = c^3\ell_0^2/(8E_\mathrm{coh})$.

\item \textbf{No graviton}: No spin-2 boson detected at any energy up to $M_\mathrm{Pl}$.

\item \textbf{Equivalence principle}: $m_i/m_g = 1$ to all orders; no deviation expected.

\item \textbf{ILG rotation curves}: $v_c^2 \propto r$ at large $r$ for isolated galaxies, 
with universal $\alpha_t \approx 0.191$.

\item \textbf{GW speed}: $c_\mathrm{GW} = c$ exactly, to all orders.
\end{enumerate}

\section{Conclusion}

RS derives all of gravitational physics from the ledger:
\begin{enumerate}
\item EFEs from action stationarity with coupling $\kappa = \pi\lambda_\mathrm{rec}^2/(\hbar c)$
\item Equivalence principle as a theorem (all masses from J-cost)
\item No graviton — gravity is emergent ledger curvature
\item Flat rotation curves from ILG with $\alpha_t = 0.5(1-\varphi^{-1})$
\item GW speed $= c$ exactly (same ledger defines both)
\end{enumerate}

\begin{thebibliography}{9}

\bibitem{RS_Axioms}
J.~Washburn, ``The Algebra of Reality: A Recognition Science Derivation
of Physical Law,'' \emph{Axioms} \textbf{15}(2), 90 (2026).
\texttt{doi:10.3390/axioms15020090}

\bibitem{gw170817} B. Abbott et al. (LIGO/Virgo), \emph{GW170817: Observation of 
Gravitational Waves from a Binary Neutron Star Inspiral}, 
Phys. Rev. Lett. 119:161101 (2017).

\bibitem{milgrom} M. Milgrom, \emph{A modification of the Newtonian dynamics as a 
possible alternative to the hidden mass hypothesis}, ApJ 270:365-370 (1983).

\bibitem{wald}
R.~M.~Wald, \emph{General Relativity}, University of Chicago Press (1984).

\end{thebibliography}

\end{document}
